%%%%%%%%%%%%%%%%%%%%%%%%%%%%%%%%%%%%%%%%%%%%%%%%%%%%%%%%%%%%%%%%%%%%%%%%%%%%%%%%
%2345678901234567890123456789012345678901234567890123456789012345678901234567890
%        1         2         3         4         5         6         7         8

\documentclass[letterpaper, 10 pt, conference]{ieeeconf}  % Comment this line out if you need a4paper

%\documentclass[a4paper, 10pt, conference]{ieeeconf}      % Use this line for a4 paper

\IEEEoverridecommandlockouts                              % This command is only needed if 
                                                          % you want to use the \thanks command

\overrideIEEEmargins                                      % Needed to meet printer 
%\usepackage{graphics} % for pdf, bitmapped graphics files
%\usepackage{epsfig} % for postscript graphics files
%\usepackage{mathptmx} % assumes new font selection scheme installed
%\usepackage{times} % assumes new font selection scheme installed
%\usepackage{amsmath} % assumes amsmath package installed
%\usepackage{amssymb}  % assumes amsmath package installed
\usepackage{cite}
\usepackage{cleveref}
\usepackage{graphicx}
\title{\LARGE \bf
Physics-Informed Deep Stereo: Active Cross-Polarization 
for Data-Efficient Transparent Obstacle Avoidance
}


\author{Po-Ting .Lin 
}


\begin{document}



\maketitle
\thispagestyle{empty}
\pagestyle{empty}


%%%%%%%%%%%%%%%%%%%%%%%%%%%%%%%%%%%%%%%%%%%%%%%%%%%%%%%%%%%%%%%%%%%%%%%%%%%%%%%%
\begin{abstract}


\end{abstract}


%%%%%%%%%%%%%%%%%%%%%%%%%%%%%%%%%%%%%%%%%%%%%%%%%%%%%%%%%%%%%%%%%%%%%%%%%%%%%%%%
\section{INTRODUCTION}

Transparent obstacles are difficult to detect using traditional methods. Their characteristics, including light transmission, spectral reflection, and refraction distortion, pose significant challenges for detection.

Traditional obstacle detection methods struggle with transparent obstacles. 
LiDAR sensors fail due to light transmission through transparent surfaces 
\cite{tibebu_lidar-based_2021}. RGB-D cameras produce incomplete depth maps with large missing regions due to spectral reflection 
and refraction distortion \cite{lysenkov_recognition_2013}. Although optical transparency does not affect 
ultrasonic sensors, they lack the spatial resolution and range required for dense 3D reconstruction in robotic navigation \cite{borenstein_vector_1991}.

Recent learning-based methods have shown promise in addressing this challenge. 
Depth completion approaches train neural networks to infer complete depth maps 
from incomplete RGB-D inputs \cite{sajjan_clear_2020}.
However, these methods are data-intensive and require tens of thousands of 
synthetic training images to achieve satisfactory performance. Active stereo methods \cite{shi_asgrasp_2024} have been proposed to improve depth estimation 
for transparent objects by leveraging structured light patterns. While these 
approaches reduce the dependency on scene texture, they still require extensive 
training data to handle the unique optical properties of transparent materials, 
such as refraction and specular reflection.

To address this limitation, we propose a physics-informed deep stereo 
framework for data-efficient transparent obstacle detection. Our approach 
leverages the physical principle that transparent objects produce highly 
polarized specular reflections. By capturing stereo pairs at orthogonal 
polarization angles (0° and 90°), the intensity difference reveals transparent 
object boundaries with high contrast, independent of texture or ambient lighting. 
This physics-based input representation enables fine-tuning of existing stereo 
networks with minimal training data, eliminating the need for large-scale 
specialized datasets.The key advantage of our approach is that it creates naturally discriminative 
patterns where transparent boundaries are clearly visible, maximizing data 
utilization efficiency—networks can learn robust features from significantly 
fewer samples because the input representation is inherently informative 
rather than ambiguous.

\section{RELATED WORK}

\subsection{Transparent Object Detection with RGB-D Sensors}

Standard RGB-D cameras fail to capture accurate depth for transparent objects 
due to light transmission and refraction. Depth completion methods aim to 
recover the corrupted depth structure caused by transparent objects through 
deep learning. These approaches can be categorized into two main strategies. 
The first, exemplified by Clear Grasp \cite{sajjan_clear_2020}, combines 
learned predictions (surface normals, masks, occlusion boundaries) with 
global optimization algorithms to reconstruct depth from over 50,000 synthetic 
training images. The second strategy, adopted by TransCG \cite{fang_transcg_2022} 
and ASGrasp \cite{shi_asgrasp_2024}, employs end-to-end networks that directly 
output completed depth maps. The key distinction between TransCG and ASGrasp 
lies in their input representation: TransCG retains the noisy depth map from 
RGB-D sensors as input, while ASGrasp discards the unreliable depth measurements 
entirely and directly utilizes RGB images paired with raw infrared (IR) stereo 
images to reconstruct geometry. Despite these different strategies, all methods 
remain data-intensive, requiring tens of thousands of training images to learn 
the complex mappings from corrupted or incomplete inputs to accurate depth.

\subsection{Polarization-Based Transparent Object Vision}

Polarization imaging reveals the geometric properties of transparent surfaces by measuring the polarization state of reflected light (degree of polarization $\rho$ and angle of polarization $\phi$), providing physical constraints for 3D reconstruction. Cui et al. \cite{cui_polarimetric_2017} explored polarimetric multi-view stereo for 3D modeling in static environments, achieving high-quality reconstruction by resolving polarization ambiguities through graph optimization. However, their multi-view acquisition (6-36 views $\times$ 7 images per view) requires several minutes, rendering it impractical for real-time applications. Zhu \& Smith \cite{zhu_depth_2019} simplified the setup to a stereo pair configuration (one polarization camera and one RGB camera), jointly solving polarization and stereo constraints to obtain metrically accurate depth maps. While this extends the applicability to broader scenarios, it still necessitates sequential capture of multiple polarization angle images and second-level optimization. Kalra et al. \cite{kalra_deep_2020} introduced deep learning as an alternative to traditional geometric optimization, employing a monocular multi-polarization camera (simultaneous capture of 4 polarization angles) to achieve real-time segmentation at 20 FPS. However, this approach requires 20,000 labeled training samples and outputs only 2D masks without depth information. All these methods rely on passive illumination, limiting their performance in low-light conditions, and exhibit inherent trade-offs among data efficiency, real-time capability, and 3D perception.

\section{METHOD}

We present a physics-informed framework for transparent obstacle avoidance that integrates polarization cues into a standard stereo vision configuration. By exploiting polarization physics to generate high-contrast observations of transparent surfaces and their boundaries, the system preserves the geometric advantages of stereo while providing more informative inputs for learning. Rather than designing a new architecture from scratch, we treat the stereo network as a reusable backbone and adapt it to the polarization-enhanced inputs via data-efficient fine-tuning. This eliminates the need for complex multi-view rigs or large-scale task-specific training datasets.

\subsection{System Overview}

\begin{figure}[h!]
    \centering
    \includegraphics[width=1\linewidth]{paper//Figures/System Optical Path Design Diagram-processed.png}
    \caption{\textbf{Schematic of the Optical Path Design.} The system utilizes a cross-polarization setup where the analyzer is orthogonal to the source polarizer. This configuration optically filters specular highlights, enhancing the visibility of transparent boundaries for the downstream stereo matching network (generated using Ansys Zemax OpticStudio).}
    \label{SOPDD-P}
\end{figure}

We employ a physics-based optical design to effectively decouple transparent object geometries from complex specular noise at the image-formation level. Specifically, we leverage the polarization-selective properties of linear polarizers. A source-side polarizer ensures that the emitted illumination maintains a fixed polarization state. Upon interaction with transparent or glossy surfaces, the specular reflection component becomes strongly polarized, whereas ambient and diffuse components remain largely unpolarized.

To exploit this physical behavior, we introduce an analyzer oriented orthogonally to the source polarizer in front of the imaging sensor. In this Cross-Polarization (CP) configuration, the strong specular noise from transparent surfaces is optically suppressed, while the unpolarized background remains visible. This mechanism effectively serves as a hardware-level attention mask, creating a physically grounded “cut-out” effect that accentuates object boundaries. Consequently, the complex inverse problem of transparency perception is simplified into a clearer representation that is easier for neural networks to interpret.

As illustrated in \cref{SOPDD-P}, our system integrates a structured-light projector with a source-side polarizer and a calibrated stereo camera pair equipped with analyzers. The outgoing light is pre-polarized, and the returning light is selectively filtered to yield high-contrast observations. Crucially, these polarization-enhanced images serve as a domain-agnostic input representation. By improving data quality at the source, we ensure that downstream stereo-matching networks can achieve high performance through data-efficient fine-tuning, without requiring specialized architectural modifications.

\subsection{Physics-Informed Input Formulation}

We decompose the illumination effects in real-world imaging into several components that contribute to the intensity observed at each pixel. Let $I(\mathbf{u})$ denote the measured image intensity at pixel location $\mathbf{u}$ This value can be attributed to three primary factors: diffuse reflection, specular reflection, and background illumination. The overall intensity at any spatial location is therefore governed by the combination of these three components.
\begin{equation}
I(\mathbf{u}) = I_{d}(\mathbf{u}) + I_{s}(\mathbf{u}) + I_{b}(\mathbf{u})
\end{equation}
To reduce the contribution of the specular reflection term in the image formation model, we introduce a pair of linear polarizing filters. According to Malus’ law, the intensity of the polarized specular component after filtering, $I_{s}^{filtered}$ ,follows $I_{s}^{filtered}\propto \cos^2\theta$ , where $\theta$ denotes the relative orientation between the polarizer and the analyzer. The specular component is maximized at $\theta=0^\circ$ and decays quadratically as the angle increases. By configuring the analyzer to be nearly orthogonal to the source-side polarizer, the specular term is strongly attenuated; in particular, at $\theta=90^\circ$ , it approaches zero because crossed polarizers effectively eliminate the polarized reflection:
\begin{equation}
I_{s}^{filtered} \approx I_{s} \cdot \cos^2(90^\circ)\approx 0
\end{equation}
Since the diffuse reflection $I_d$ and background illumination $I_b$ are largely unpolarized, they are transmitted through the analyzer with a constant attenuation factor $\alpha$ (typically $\alpha \approx 0.5$ for ideal polarizers). Consequently, the final captured image $I^{CP}$ can be formulated as:
\begin{equation}
I^{CP} \approx \alpha \cdot (I_{d}(\mathbf{u}) + I_{b}(\mathbf{u}))
\end{equation}
This selective suppression significantly reduces the specular contribution while leaving the diffuse and background components largely unaffected, thereby enhancing boundary contrast and providing more discriminative cues for downstream learning.

\subsection{Standardized Backbone Utilization}

\subsection{Optimization and Generalization Objectives}

\section{EXPERIMENTS}
\section{RESULTS}
\section{CONCLUSIONS}
\section*{APPENDIX}
\bibliographystyle{IEEEtran.bst}  % 指定使用 IEEE 標準格式
\bibliography{IEEEfull.bib}










\end{document}
